%% kapa: exam=1 round=27 year=2016 period=1 form=A subject=real_estate subject_ko=부동산학원론 seq=32 number=112
%% source: Q-Net 감정평가사 기출문제 (cst00309) artlSeq=5169402 file=extracted/1008/5169402/제27회 감정평가사 1차 문제/제27회 감정평가사 1교시문제(A형).pdf
%% conversion: kordoc 4.24.51 -> tools/gen1.py (auto); review=proofread-sonnet; flags=none
\begin{question}{112}
주택법령상 주택의 정의에 관한 설명으로 옳지 않은 것은?
\begin{choices}
  \item 주택은 세대의 구성원이 장기간 독립된 주거생활을 할 수 있는 구조로 된 건축물의 전부 또는 일부 및 그 부속토지를 말한다.
  \item 준주택은 주택 외의 건축물과 그 부속토지로서 주거시설로 이용가능한 시설 등을 말한다.
  \item 공동주택은 건축물의 벽ㆍ복도ㆍ계단이나 그 밖의 설비 등의 전부 또는 일부를 공동으로 사용하는 각 세대가 하나의 건축물 안에서 각각 독립된 주거생활을 할 수 있는 구조로 된 주택을 말한다.
  \item 민영주택은 국민주택등을 제외한 주택을 말한다.
  \item 세대구분형 공동주택은 300세대 미만의 국민주택규모에 해당하는 주택으로서 단지형 연립주택, 단지형 다세대주택, 원룸형 주택으로 분류한다.
\end{choices}
\end{question}
\answer{5}
